%======================================================================
\chapter{Introduction}\label{ch:intro}
%======================================================================

\section{What this volume is about}

This is the last volume of the notes. It proves the classification of homogeneous plane continua (\vone{Theorem~0.1}):

\begin{theorem}[Theorem~\ref{thm:classification}]
A nondegenerate plane continuum is homogeneous if and only if it is a simple closed curve, a pseudo-arc, or a circle of pseudo-arcs.
\end{theorem}

and its plane form (Theorem~\ref{thm:classification-plane}): there are three plane continua, a simple closed curve, a pseudo-arc and a circle of
pseudo-arcs, each homogeneous, no two homeomorphic, and every nondegenerate homogeneous plane continuum is homeomorphic to exactly one of them.

Volume~6 settled the indecomposable case: a nondegenerate homogeneous indecomposable plane continuum is a pseudo-arc (\vsix{Theorem~7.4}). Volume~7
settled most of the decomposable case: a decomposable homogeneous plane continuum $M$ is a simple closed curve, or there is a continuous, open, surjective
map $f\colon M\to\Sph$ whose fibres are pseudo-arcs and are terminal in $M$ (\vseven{Theorem~9.8}). What is missing is the theory of the third class.

A \emph{circle of pseudo-arcs}, in the sense of Bing and Jones (1959), is a circle-like continuum $X$ with a continuous, open, surjective map $X\to\Sph$
whose fibres are pseudo-arcs. Three things must be proved about it.
\begin{enumerate}
\item \textbf{Uniqueness and homogeneity.} Any two circles of pseudo-arcs are homeomorphic, and each is homogeneous. Bing and Jones derive both from their
Theorem~10: between two ``continuous snake-like arcs of pseudo-arcs'' over $\I$, any homeomorphisms of the end fibres extend to a homeomorphism that
carries each fibre onto the fibre over the same point.
\item \textbf{Recognition.} The map of Volume~7 makes $M$ a circle of pseudo-arcs. Volume~7 does not show that $M$ is circle-like, and an open map onto
$\Sph$ with pseudo-arc fibres does not force it by itself. The map of Volume~7 is \emph{atomic}, and with atomicity the halves of $M$ over the two
half-circles are snake-like. Theorem~10 then makes $M$ homeomorphic to any circle of pseudo-arcs, hence circle-like (Domokos and Prajs, Theorem~6.1).
\item \textbf{Existence in the plane.} Some plane continuum is a circle of pseudo-arcs. Bing and Jones assert this (their \S7) but do not construct the
example. These notes build one by a different construction, a circular analogue of the plane pseudo-arc of \vone{Chapter~5}.
\end{enumerate}

\section{The chapters}

\begin{itemize}
\item \textbf{Chapter~\ref{ch:circles}: Arcs and circles of pseudo-arcs.} The definitions, the two halves of a map onto the circle, and three facts that
distinguish the circle, the pseudo-arc and circles of pseudo-arcs.
\item \textbf{Chapter~\ref{ch:atomic}: Atomic arcs are snake-like.} Bing's proof of his 1951 Theorem~11, with the fibres of an atomic map in place of his
layers: an atomic map onto $\I$ with snake-like hereditarily indecomposable fibres has a snake-like domain. Hence the halves of an atomic map are
snake-like.
\item \textbf{Chapter~\ref{ch:chains}: Chains in the sense of Bing and Jones.} Chains whose nonadjacent links are a positive distance apart; chain maps;
private points; runs; folding a chain along a walk; the form of Bing and Jones's Theorem~6 that is used.
\item \textbf{Chapter~\ref{ch:arcs}: Snake-like arcs of continua.} Bing and Jones's Theorems~3, 4 and~5 and Step~1(b) of Theorem~10. Theorem~5 is false
as printed; it is proved in a repaired form.
\item \textbf{Chapter~\ref{ch:thm78}: Theorems 7 and 8.} Chains following a chain map, on one pseudo-arc (Theorem~7) and on a whole arc of pseudo-arcs
(Theorem~8 and its corollary).
\item \textbf{Chapter~\ref{ch:zigzag}: The homeomorphism from a zigzag.} Two sequences of chains with chain maps going back and forth define a
homeomorphism (Bing and Jones, pp.~188--189).
\item \textbf{Chapter~\ref{ch:thm10}: Theorem 10.} One step of the construction, made alternately on the two sides, and the induction.
\item \textbf{Chapter~\ref{ch:uniqueness}: Uniqueness, homogeneity and recognition.} Gluing two half homeomorphisms; Domokos and Prajs, Theorem~6.1.
\item \textbf{Chapter~\ref{ch:circdescr}: Circular descriptions.} Nested circular chains whose links carry arcs of the base circle; a recognition theorem
for the continuum they describe.
\item \textbf{Chapter~\ref{ch:sweep}: Sweep patterns.} Explicit integer sequences satisfying the conditions of Chapter~\ref{ch:circdescr}, and the
combinatorial ``sweep lemma'' that makes them crooked.
\item \textbf{Chapter~\ref{ch:plane}: A circle of pseudo-arcs in the plane.} The circular analogue of the planar step of \vone{Proposition~5.14}, by a
shear on the universal cover of the punctured plane, and the assembly of a circular description in $\R^2$.
\item \textbf{Chapter~\ref{ch:classification}: The classification theorem.}
\end{itemize}

The main results of the volume are:

\begin{theorem}[Bing 1951, for atomic maps; Theorem~\ref{thm:atomic-snake}]
Let $X$ be a compact metric space and $f\colon X\to\I$ a continuous atomic map whose fibres are snake-like and hereditarily indecomposable. Then $X$ is
snake-like.
\end{theorem}

\begin{theorem}[Bing and Jones, Theorem~10; Theorem~\ref{thm:bj10}]
Let $f\colon X\to\I$ and $g\colon Y\to\I$ be continuous, open and surjective, with pseudo-arc fibres, where $X$ and $Y$ are snake-like compact metric spaces.
Let $h_0\colon\fib f0\to\fib g0$ and $h_1\colon\fib f1\to\fib g1$ be homeomorphisms. Then there is a homeomorphism $H\colon X\to Y$ with $g\circ H=f$ that
extends $h_0$ and $h_1$.
\end{theorem}

\begin{theorem}[Theorems~\ref{thm:unique}, \ref{thm:homog}, \ref{thm:dp61}]
Any two circles of pseudo-arcs are homeomorphic, and each is homogeneous. A compact metric space with an open atomic map onto $\Sph$ whose fibres are
pseudo-arcs is a circle of pseudo-arcs.
\end{theorem}

\begin{theorem}[Theorem~\ref{thm:plane-example}]
Some plane continuum is a circle of pseudo-arcs.
\end{theorem}

\section{Where the formal proof differs from the literature}\label{sec:intro-diff}

\begin{enumerate}
\item \textbf{Halves of an atomic map.} Domokos and Prajs obtain the snake-likeness of the halves from a proposition of Hagopian, Marsh and Prajs on
``atriodic folders''. The formal proof uses instead Bing's 1951 argument (Chapter~\ref{ch:atomic}); no atriodicity is needed.
\item \textbf{Chains.} Bing and Jones require nonadjacent links of a chain to be a positive distance apart; the earlier volumes only ask them to be disjoint.
The theorems are stated with the weaker hypothesis and the stronger conclusion wherever possible, and a shrinking lemma passes between the two
(Theorem~\ref{thm:shrink}).
\item \textbf{Theorems 3, 5, 7 and 8.} The chain printed in the proof of Theorem~3 is not a chain; it is repaired at its two junctions. Theorem~5 is false
as printed (Proposition~\ref{prop:bj5-false}); a repaired form is proved and used, on a window of the arc. The second sentence of Theorem~7 is false for
chains of one link and its printed proof fails for three links; it is proved for every chain of at least three links with a different auxiliary chain. In
Theorem~8 the printed links need not form a chain; they are restricted to the fibres over their intervals.
\item \textbf{Theorem 10.} Bing and Jones extend the homeomorphism to cut fibres in the interior of the arc at each step, using their Theorem~9. The formal
proof uses the homeomorphism only on the two end fibres, and passes from one sub-arc to the next through overlapping windows. Theorem~9 is therefore not
used. One step serves for both sides of the induction, and the induction starts from chains of one link.
\item \textbf{The passage to circles.} ``It follows from Theorem~10 that any two circles of pseudo-arcs are topologically equivalent and that each is
homogeneous'' is not argued by Bing and Jones; it is supplied by cutting the circle into two halves (Chapter~\ref{ch:uniqueness}).
\item \textbf{The plane example.} Bing and Jones's \S7 asserts a sequence of circular chains with nine properties without constructing it, and does not argue
that the resulting continuum has pseudo-arc fibres; their condition~(9) does not visibly give it. The example of Chapters~\ref{ch:circdescr}--\ref{ch:plane}
is built by another route.
\end{enumerate}

\section{Prerequisites and conventions}\label{sec:intro-conv}

The reader should know chains, crookedness, descriptions and the plane pseudo-arc (Volume~1); pattern following, the uniqueness and homogeneity of the
pseudo-arc (Volume~2); boundary bumping, snake-like and circle-like continua (\vthree{Chapters~1 and~5}); and the main theorems of Volumes~6 and~7.

\begin{itemize}
\item \textbf{Continua.} A continuum is compact, connected and metric; it may be a point unless ``nondegenerate'' is said. A pseudo-arc is a nondegenerate,
snake-like, hereditarily indecomposable continuum (\vone{Definition~2.14}).
\item \textbf{Chains} are as in \vone{Definition~2.1}: open links $D_0,\dots,D_{n-1}$ with $D_i\cap D_j\ne\emptyset$ if and only if $\abs{i-j}\le1$. Links are
numbered from~$0$. Bing and Jones number from~$1$; quotations are translated.
\item \textbf{Fibres.} For $f\colon X\to\I$ and $x\in\I$ the fibre $\fib fx$ is also written $p_x$ (and $q_x$ for $g\colon Y\to\I$), following Bing and Jones.
For $a\le b$ we write $\G ab=\fib f{[a,b]}$.
\item \textbf{The circle} is $\Sph=\set{z\in\C:\abs z=1}$, and $e^{i\theta}$ is written $\exp(i\theta)$. The plane carries the sup metric, as in the earlier
volumes.
\end{itemize}

\paragraph{Formal counterparts.} As in the earlier volumes, many results are followed by a line ``Formal counterpart:'' naming the corresponding declaration of the
Lean development (all names begin with \lean{HPC.}). These lines may be skipped. Appendix~\ref{app:formal} collects them.

\paragraph{Proofs in outline.} Some long constructions (Bing's surgery and Fugate's lemma in Chapter~\ref{ch:atomic}, Theorem~7 in Chapter~\ref{ch:thm78}, the
step of Theorem~10, the planar step of Chapter~\ref{ch:plane}) are given as proofs in outline, marked as such. The complete arguments are in the declarations named in the ``Formal counterpart'' lines.

\paragraph{Exercises.} Each chapter ends with exercises. None of them is used later.
